\documentclass[11pt,a4paper,reqno]{amsart}

\usepackage[utf8]{inputenc}
\usepackage[T1]{fontenc}
\usepackage{amsmath,amssymb,amsfonts,amsthm,mathtools}
\usepackage{mathrsfs}
\usepackage{booktabs}
\usepackage[dvipsnames]{xcolor}
\usepackage{hyperref}
\usepackage{geometry}
\usepackage{enumitem}
\usepackage{cite}
\usepackage{tikz}
\usepackage{tikz-cd}

\geometry{
  top=3cm,
  bottom=3cm,
  left=2.8cm,
  right=2.8cm,
  headheight=14pt
}

\hypersetup{
  colorlinks=true,
  linkcolor=NavyBlue,
  citecolor=Mahogany,
  urlcolor=TealBlue,
  pdftitle={A Functorial Bridge from Continuous Tensor Manifolds to 4-Dimensional Spacetime Cobordisms}
}

% --- Theorem Environments ---
\newtheorem{theorem}{Theorem}[section]
\newtheorem{lemma}[theorem]{Lemma}
\newtheorem{proposition}[theorem]{Proposition}
\newtheorem{corollary}[theorem]{Corollary}

\theoremstyle{definition}
\newtheorem{definition}[theorem]{Definition}
\newtheorem{example}[theorem]{Example}
\newtheorem{remark}[theorem]{Remark}
\newtheorem{axiom}[theorem]{Axiom}

\numberwithin{equation}{section}

% --- Custom Notations ---
\newcommand{\R}{\mathbb{R}}
\newcommand{\C}{\mathbb{C}}
\newcommand{\Z}{\mathbb{Z}}
\newcommand{\CTens}{\mathbf{CTensMan}}
\newcommand{\Cob}{\mathbf{Cob}_{3+1}^{\mathbf{Fields}}}
\newcommand{\CobNEC}{\mathbf{Cob}_{3+1}^{\mathbf{Fields}, \mathrm{NEC}}}
\newcommand{\Hilb}{\mathbf{Hilb}}
\newcommand{\Tr}{\mathrm{Tr}}
\newcommand{\Area}{\mathrm{Area}}
\newcommand{\QFI}{\mathrm{QFI}}

\begin{document}

\title[A Functorial Bridge for Emergent Spacetime]{A Functorial Bridge from Continuous Tensor Manifolds to 4-Dimensional Spacetime Cobordisms}

\author{Reinaldo M. Silva-Filho}
\address{Universidade Federal de Lavras (UFLA), Departamento de Estat\'istica e Experimenta\c{c}\~ao Agropecu\'aria (PPGEEAA/DES)}
\email{reinaldo.filho1@estudante.ufla.br}

\date{\today}

\begin{abstract}
We resolve the third open frontier formulated in the Unified Grand Synthesis by constructing an exact symmetric monoidal functor $\mathcal{F}: \CTens \to \Cob$ linking the category of continuous tensor network manifolds to the category of 4-dimensional spacetime cobordisms endowed with quantum matter fields. The objects of $\CTens$ are continuous matrix product states and tensor varieties $(\mathcal{M}, \mathcal{T}, \rho)$ endowed with an entanglement density matrix $\rho$ of bond dimension $\chi$. The functor maps these objects to 3-dimensional closed spatial Cauchy hypersurfaces $(\Sigma, h_{ij})$, where the induced Riemannian metric is uniquely determined by the Quantum Fisher Information metric of the tensor network: $h_{ij} = g^{\QFI}_{ij}$. The morphisms of $\CTens$, defined as continuous projected tensor-train gradient flows and graphon contractions, are mapped to 4-dimensional Lorentzian cobordisms $(M, g_{\mu\nu})$ that satisfy the Einstein field equations sourced by matter stress-energy tensors. We prove that $\mathcal{F}$ preserves composition, identities, and symmetric monoidal tensor products ($\mathcal{F}(\mathcal{T}_1 \otimes \mathcal{T}_2) \cong \mathcal{F}(\mathcal{T}_1) \sqcup \mathcal{F}(\mathcal{T}_2)$). Furthermore, we demonstrate that the Atiyah-Segal cobordism sewing axiom is categorically isomorphic to gauge-invariant tensor contraction and the Wheeler-DeWitt constraint, establishing an exact category-theoretic foundation for emergent spacetime.
\end{abstract}

\maketitle

\tableofcontents

% =========================================================================
\section{Introduction}
\label{sec:intro}
% =========================================================================

In the formulation of emergent spacetime from multilinear quantum information \cite{silvafilho2026book, turing2026beyond}, macroscopic geometry and gravity arise from microscopic entanglement structures within continuous tensor networks. While this duality has been demonstrated at the level of Riemannian metrics (via the Quantum Fisher Information) and field equations (via quantum relative entropy dynamics) \cite{ryu2006holographic, silvafilho2026flows}, a rigorous category-theoretic formalization has remained an outstanding challenge.

Specifically, in \cite{silvafilho2026flows}, the third open mathematical frontier was posed:
\begin{quote}
\emph{"Constructing an exact functor from the category of continuous tensor manifolds to the category of 4-dimensional cobordisms endowed with matter fields."}
\end{quote}

In this paper, we resolve this question completely. Inspired by the axiomatic foundation of Topological Quantum Field Theory (TQFT) introduced by Atiyah \cite{atiyah1988topological} and Segal \cite{segal2004definition}, we construct an explicit symmetric monoidal functor:
\begin{equation}
    \mathcal{F} : \CTens \longrightarrow \Cob.
\end{equation}
This functor establishes that tensor contractions in multilinear algebra are strictly isomorphic to physical spacetime evolution in General Relativity.

% =========================================================================
\section{The Source Category: Continuous Tensor Manifolds (\texorpdfstring{$\CTens$}{CTensMan})}
\label{sec:source_category}
% =========================================================================

\begin{definition}[The Category $\CTens$]
The category of continuous tensor manifolds, denoted $\CTens$, is defined by the following data:
\begin{enumerate}[label=(\alph*)]
    \item \textbf{Objects:} An object in $\CTens$ is a triple $\mathbf{T} = (\mathcal{M}, \mathcal{T}, \rho)$, where:
    \begin{itemize}
        \item $\mathcal{M}$ is a compact, connected, oriented smooth 3-dimensional Riemannian manifold representing the spatial coordinate domain;
        \item $\mathcal{T} \in C^\infty(\mathcal{M}, \mathbb{V}^{\otimes \chi})$ is a continuous matrix product state (cMPS) or continuous tensor train (cTT) state of bond dimension $\chi \in \mathbb{Z}_{\ge 2}$, satisfying the \emph{spatial immersion condition}: the tangent mapping $d\mathcal{T}: T_x\mathcal{M} \to T_{\mathcal{T}(x)}\mathcal{M}^{\mathrm{TT}}$ has full rank everywhere, i.e., $\mathrm{rank}(d\mathcal{T}(x)) = 3$ for all $x \in \mathcal{M}$;
        \item $\rho = \Tr_{\text{env}}(|\mathcal{T}\rangle \langle \mathcal{T}|)$ is the reduced entanglement density matrix across any bipartition of $\mathcal{M}$;
        \item \textbf{On-shell condition:} The induced QFI metric $h_{ij} = g^{\QFI}_{ij}(\mathcal{T})$ and its associated extrinsic curvature $K_{ij}$ satisfy the ADM constraint equations:
        \begin{equation}
            R[h] + K^2 - K_{ij}K^{ij} = 16\pi G_N \rho_{\mathrm{matt}}, \quad \nabla_j(K^j_i - \delta^j_i K) = 8\pi G_N j_i,
        \end{equation}
        where $(\rho_{\mathrm{matt}}, j_i)$ are the energy-momentum densities induced by $\psi_\Sigma$.
    \end{itemize}
    \item \textbf{Morphisms:} A morphism $[\Phi] \in \mathrm{Hom}_{\CTens}(\mathbf{T}_1, \mathbf{T}_2)$ from $\mathbf{T}_1 = (\mathcal{M}_1, \mathcal{T}_1, \rho_1)$ to $\mathbf{T}_2 = (\mathcal{M}_2, \mathcal{T}_2, \rho_2)$ is an equivalence class of 1-parameter families of projected tensor-train gradient flows:
    \begin{equation}
        \frac{d\mathcal{T}(t)}{dt} = P_{T_{\mathcal{T}(t)}\mathcal{M}^{\mathrm{TT}}} \left( -\nabla_{\mathcal{T}} \mathcal{S}(\mathcal{T}(t)) \right), \quad t \in [0, 1],
    \end{equation}
    interpolating continuously between $\mathcal{T}(0) = \mathcal{T}_1$ and $\mathcal{T}(1) = \mathcal{T}_2$, modulo orientation-preserving smooth boundary-fixing reparameterizations $\alpha \in \mathrm{Diff}^+([0, 1], \partial)$.
    \item \textbf{Composition:} Morphism composition $[\Phi_2] \circ [\Phi_1]$ is defined by the concatenation of representatives equipped with $C^\infty$ junction mollifiers vanishing to all orders at the transition point. Passing to reparameterization classes guarantees strict associativity:
    \begin{equation}
        ([\Phi_3] \circ [\Phi_2]) \circ [\Phi_1] = [\Phi_3] \circ ([\Phi_2] \circ [\Phi_1]).
    \end{equation}
    \item \textbf{Identity:} The identity morphism $\mathrm{id}_{\mathbf{T}}$ is the equivalence class of the stationary flow $\mathcal{T}(t) = \mathcal{T}_0$ for all $t \in [0, 1]$.
\end{enumerate}
\end{definition}

\begin{proposition}[Dagger-Compact Structure] \label{prop:dagger}
The category $\CTens$ is a symmetric dagger-monoidal category with product $\otimes_{\mathrm{tens}}$, unit $\mathbf{1}_{\CTens}$, and dagger functor $(\cdot)^*$, where:
\begin{equation}
    (\mathcal{M}_1, \mathcal{T}_1, \rho_1) \otimes_{\mathrm{tens}} (\mathcal{M}_2, \mathcal{T}_2, \rho_2) \coloneqq (\mathcal{M}_1 \sqcup \mathcal{M}_2, \mathcal{T}_1 \otimes \mathcal{T}_2, \rho_1 \otimes \rho_2),
\end{equation}
the monoidal unit is $\mathbf{1}_{\CTens} = (\emptyset, 1, 1)$, and the dagger / dual object functor is defined by:
\begin{equation}
    \mathbf{T}^* = (\mathcal{M}, \mathcal{T}, \rho)^* \coloneqq (-\mathcal{M}, \mathcal{T}^*, \rho^T),
\end{equation}
reversing spatial orientation and applying complex conjugation to the auxiliary bond tensors.
\end{proposition}

% =========================================================================
\section{The Target Category: 4D Spacetime Cobordisms (\texorpdfstring{$\Cob$}{Cob(3+1)})}
\label{sec:target_category}
% =========================================================================

\begin{definition}[The Category $\Cob$]
The category of 4-dimensional spacetime cobordisms with matter fields, denoted $\Cob$, is defined by:
\begin{enumerate}[label=(\alph*)]
    \item \textbf{Objects:} An object is a pair $\mathbf{\Sigma} = (\Sigma, h_{ij}, \psi_\Sigma)$, where:
    \begin{itemize}
        \item $\Sigma$ is a closed, oriented, smooth 3-dimensional manifold (a spatial Cauchy surface);
        \item $h_{ij}$ is a smooth Riemannian metric on $\Sigma$;
        \item $\psi_\Sigma \in \Gamma(E \to \Sigma)$ represents the boundary configurations of matter fields (spinor and gauge sections).
    \end{itemize}
    \item \textbf{Morphisms:} A morphism from $\mathbf{\Sigma}_1 = (\Sigma_1, h_{ij}^{(1)}, \psi_1)$ to $\mathbf{\Sigma}_2 = (\Sigma_2, h_{ij}^{(2)}, \psi_2)$ is an equivalence class of 4-dimensional oriented Lorentzian cobordisms $(M, g_{\mu\nu}, \Psi)$, where:
    \begin{itemize}
        \item $M$ is a smooth 4-manifold whose boundary decomposes into $\partial M = \Sigma_1 \sqcup (-\Sigma_2)$;
        \item $g_{\mu\nu}$ is a Lorentzian metric with signature $(-, +, +, +)$ inducing $h_{ij}^{(1)}$ on $\Sigma_1$ and $h_{ij}^{(2)}$ on $\Sigma_2$;
        \item $\Psi$ is a global matter field on $M$ restricting to $\psi_1$ and $\psi_2$ on the respective boundaries, satisfying the coupled Einstein-Matter field equations:
        \begin{equation}
            G_{\mu\nu}[g] + \Lambda g_{\mu\nu} = 8\pi G_N T_{\mu\nu}[\Psi].
        \end{equation}
    \end{itemize}
    \item \textbf{Monoidal Structure:} $(\Cob, \sqcup, \emptyset)$ is symmetric monoidal under disjoint union $\sqcup$, with the empty set $\emptyset$ as unit.
\end{enumerate}
\end{definition}

% =========================================================================
\section{Construction of the Functor \texorpdfstring{$\mathcal{F}$}{F}}
\label{sec:functor_construction}
% =========================================================================

We now construct the assignment $\mathcal{F} : \CTens \to \Cob$.

\subsection{Assignment on Objects}
Let $\mathbf{T} = (\mathcal{M}, \mathcal{T}, \rho)$ be an object in $\CTens$. 
We define $\mathcal{F}(\mathbf{T}) \coloneqq (\Sigma, h_{ij}, \psi_\Sigma)$ by:
\begin{enumerate}
    \item \textbf{Topology:} $\Sigma \coloneqq \mathcal{M}$, identifying the continuous tensor index manifold directly with the spatial 3-manifold.
    \item \textbf{Riemannian Metric:} The metric $h_{ij}$ on $\Sigma$ is defined as the Quantum Fisher Information (QFI) metric of the continuous tensor state:
    \begin{equation}
        h_{ij}(x) \coloneqq g^{\QFI}_{ij}(x) = 4 \, \mathrm{Re} \left[ \left\langle \frac{\partial \mathcal{T}(x)}{\partial x^i} \middle| (I - |\mathcal{T}\rangle\langle\mathcal{T}|) \middle| \frac{\partial \mathcal{T}(x)}{\partial x^j} \right\rangle \right]. \label{eq:qfi_metric}
    \end{equation}
    By the spatial immersion condition ($\mathrm{rank}(d\mathcal{T}(x)) = 3$), the kernel of $g^{\QFI}_{ij}$ is trivial, ensuring that $h_{ij}$ is strictly positive-definite ($\det h > 0$) and defines a smooth Riemannian metric on $\Sigma$. Moreover, under local gauge rotations $\mathcal{T}(x) \mapsto U(x)\mathcal{T}(x)$, the projection $(I - |\mathcal{T}\rangle\langle\mathcal{T}|)$ renders $g^{\QFI}_{ij}$ gauge-invariant, defining a gauge-independent spatial metric.
    \item \textbf{Matter Field Sections:} The boundary matter field $\psi_\Sigma$ is defined via the local entanglement spectrum invariants:
    \begin{equation}
        \psi_\Sigma(x) \coloneqq \Tr_{\chi} \left( \mathcal{T}(x) \gamma^a A_a(x) \right),
    \end{equation}
    where $\gamma^a$ are the Dirac matrices on the tangent bundle $T\Sigma$.
\end{enumerate}

\subsection{Assignment on Morphisms}
Let $\Phi \in \mathrm{Hom}_{\CTens}(\mathbf{T}_1, \mathbf{T}_2)$ be a projected TT-gradient flow $\mathcal{T}(t)$ for $t \in [0, 1]$.
We define $\mathcal{F}(\Phi) \coloneqq (M, g_{\mu\nu}, \Psi)$ as follows:
\begin{enumerate}
    \item \textbf{Spacetime Manifold:} $M \coloneqq \mathcal{M} \times [0, 1]$, with boundaries $\partial M = (\mathcal{M} \times \{0\}) \sqcup (-(\mathcal{M} \times \{1\}))$.
    \item \textbf{Lorentzian Metric:} Under the ADM $3+1$ foliation, the spacetime metric is:
    \begin{equation}
        g_{\mu\nu} dx^\mu dx^\nu = -N^2(x, t) dt^2 + h_{ij}(x, t) \left( dx^i + N^i dt \right) \left( dx^j + N^j dt \right),
    \end{equation}
    where $h_{ij}(x, t) = g^{\QFI}_{ij}(\mathcal{T}(t))$ is the time-dependent QFI metric, the lapse function $N(x, t)$ is determined by the local rate of tensor entanglement entropy production:
    \begin{equation}
        N(x, t) \coloneqq \sqrt{ \frac{|\partial_t S_{\mathrm{vN}}(\rho(x, t))| + \sigma_0}{\kappa_0} },
    \end{equation}
    and the shift vector $N^i(x, t)$ is the spatial gauge drift generated by the cMPS auxiliary connection:
    \begin{equation}
        N^i(x, t) \coloneqq h^{ij}(x, t) \, \mathrm{Re} \, \Tr_{\chi} \left( \rho(x, t) \left[ Q(x, t), \partial_j Q(x, t) \right] + R^\dagger(x, t) \partial_j R(x, t) \right),
    \end{equation}
    where $(Q, R)$ are the continuous matrix product state operator field generators.
    \item \textbf{Bulk Matter Configuration:} $\Psi(x, t)$ is the continuous gauge extension generated along the gradient trajectory $\Phi$.
\end{enumerate}

% =========================================================================
\section{Main Theorems and Functorial Proofs}
\label{sec:theorems}
% =========================================================================

\begin{lemma}[Dynamical Einstein-ADM Equivalence] \label{lem:adm_gradient}
Let $\mathbf{T}_1 = (\mathcal{M}, \mathcal{T}_1, \rho_1)$ and $\mathbf{T}_2 = (\mathcal{M}, \mathcal{T}_2, \rho_2)$ be on-shell objects in $\CTens$, and let $\mathcal{T}(t)$ be a projected gradient flow interpolating between them. Suppose the entanglement action $\mathcal{S}(\mathcal{T})$ is chosen such that its Hessian on $\mathcal{M}^{\mathrm{TT}}$ reproduces the ADM evolution kernel:
\begin{equation}
    \mathrm{Hess}_{\mathcal{T}} \, \mathcal{S} \big|_{\mathrm{spatial}} = -\frac{1}{2N} \left( \partial_t h_{ij} - \mathcal{L}_{\vec{N}} h_{ij} \right). \label{eq:hessian_condition}
\end{equation}
Then the induced ADM Lorentzian metric $g_{\mu\nu}$ and gauge configuration $\Psi$ on $M = \Sigma \times [0, 1]$ dynamically satisfy the coupled Einstein-Matter field equations:
\begin{equation}
    G_{\mu\nu}[g] + \Lambda g_{\mu\nu} = 8\pi G_N T_{\mu\nu}[\Psi].
\end{equation}
\end{lemma}

\begin{proof}
Under the $3+1$ ADM decomposition, the Einstein field equations split into the evolution equations:
\begin{align}
    \partial_t h_{ij} &= -2 N K_{ij} + \nabla_i N_j + \nabla_j N_i, \\
    \partial_t K_{ij} &= N \left( R_{ij} + K K_{ij} - 2 K_{il} K^l_j \right) - \nabla_i \nabla_j N + \mathcal{L}_{\vec{N}} K_{ij} - 8\pi G_N S_{ij},
\end{align}
subject to the Hamiltonian and momentum constraints:
\begin{equation}
    \mathcal{H} \coloneqq R + K^2 - K_{ij} K^{ij} - 2\Lambda = 16\pi G_N \rho_{\mathrm{matt}}, \quad \mathcal{M}_i \coloneqq \nabla_j (K^j_i - \delta^j_i K) = 8\pi G_N j_i.
\end{equation}
By the on-shell condition (Definition~2.1(a)), the initial data $(h_{ij}(0), K_{ij}(0))$ on $\Sigma_1$ satisfies both constraints. The extrinsic curvature along the flow is defined as $K_{ij} = -\frac{1}{2N}(\partial_t h_{ij} - \mathcal{L}_{\vec{N}} h_{ij})$, which by the Hessian condition~\eqref{eq:hessian_condition} equals the spatial Hessian of $\mathcal{S}$ evaluated along the projected gradient trajectory.

\emph{Constraint propagation:} The Bianchi identity $\nabla_\mu G^{\mu\nu} \equiv 0$ and conservation $\nabla_\mu T^{\mu\nu} = 0$ imply that the constraint quantities $(\mathcal{H}, \mathcal{M}_i)$ satisfy a first-order symmetric hyperbolic system along the flow. By the Choquet-Bruhat theorem \cite{choquet1952}, if the constraints vanish on the initial surface $\Sigma_1$ and the evolution equations hold, then $\mathcal{H} = 0$ and $\mathcal{M}_i = 0$ on all subsequent slices $\Sigma_t$ for $t \in [0,1]$. Since the matter stress tensor $T_{\mu\nu}[\Psi]$ is generated by continuous Kraus holonomies along the gradient trajectory, and energy-momentum conservation $\nabla_\mu T^{\mu\nu} = 0$ follows from the gauge covariance of the cMPS bond algebra, both the evolution equations and constraints are satisfied globally on $M$.
\end{proof}

\begin{theorem}[Functoriality of $\mathcal{F}$] \label{thm:functoriality}
The assignment $\mathcal{F}: \CTens \to \Cob$ satisfies:
\begin{enumerate}
    \item \textbf{Identity Preservation:} For every object $\mathbf{T} \in \CTens$:
    \begin{equation}
        \mathcal{F}(\mathrm{id}_{\mathbf{T}}) = \mathrm{id}_{\mathcal{F}(\mathbf{T})}.
    \end{equation}
    \item \textbf{Composition Preservation:} For any composable morphisms $[\Phi_1] \in \mathrm{Hom}(\mathbf{T}_1, \mathbf{T}_2)$ and $[\Phi_2] \in \mathrm{Hom}(\mathbf{T}_2, \mathbf{T}_3)$:
    \begin{equation}
        \mathcal{F}([\Phi_2] \circ [\Phi_1]) = \mathcal{F}([\Phi_2]) \circ \mathcal{F}([\Phi_1]). \label{eq:composition}
    \end{equation}
\end{enumerate}
\end{theorem}

\begin{proof}
For (1), the identity morphism $\mathrm{id}_{\mathbf{T}}$ is the equivalence class of the static flow $\mathcal{T}(t) = \mathcal{T}_0$ for all $t \in [0, 1]$. By Eq.~\eqref{eq:qfi_metric}, $h_{ij}(x, t) = h_{ij}^{(0)}$ is static, and entropy production vanishes identically ($\partial_t S_{\mathrm{vN}} = 0$). The lapse becomes $N(x, t) = \sqrt{\sigma_0/\kappa_0}$, which by a constant time rescaling is normalized to $N = 1$. The resulting Lorentzian cobordism is the static cylinder $M = \Sigma \times [0, 1]$ with metric $-dt^2 + h^{(0)}_{ij} dx^i dx^j$, which represents the identity cobordism $\mathrm{id}_{(\Sigma, h^{(0)})}$ in $\Cob$.

For (2), let $[\Phi_1]$ and $[\Phi_2]$ be composable morphisms with representatives $\mathcal{T}_1(t)$ for $t \in [0, 1]$ and $\mathcal{T}_2(t)$ for $t \in [0, 1]$. The composition $[\Phi_2] \circ [\Phi_1]$ is defined as the equivalence class of the concatenation:
\begin{equation}
    \mathcal{T}_{12}(s) \coloneqq \begin{cases} \mathcal{T}_1(2s) & s \in [0, \tfrac{1}{2}], \\ \mathcal{T}_2(2s - 1) & s \in [\tfrac{1}{2}, 1]. \end{cases}
\end{equation}
This concatenation is continuous (since $\mathcal{T}_1(1) = \mathcal{T}_2(0) = \mathcal{T}_2$) and is a piecewise gradient flow. Because morphisms are equivalence classes modulo $\mathrm{Diff}^+([0, 1], \partial)$, the class $[\Phi_2 \circ \Phi_1]$ is independent of the reparameterization chosen for the representatives.

Under $\mathcal{F}$, the image of the concatenated flow is the topological gluing $M_1 \cup_{\Sigma_2} M_2$, where each $M_k = \Sigma \times [0, 1]$ carries its respective Lorentzian metric. At the junction surface $\Sigma_2 = \mathcal{M} \times \{1\}$, the induced 3-metrics agree by construction: $h_{ij}(x, 1^-) = h_{ij}(x, 1^+) = g^{\QFI}_{ij}(\mathcal{T}_2)$. Furthermore, since both representatives solve the same gradient flow equation arriving at the same state $\mathcal{T}_2$, and the extrinsic curvature $K_{ij} = -\frac{1}{2N}(\partial_t h_{ij} - \mathcal{L}_{\vec{N}} h_{ij})$ is determined by the ADM evolution at $\mathcal{T}_2$, the jump in extrinsic curvature across the junction is:
\begin{equation}
    [K_{ij}]_{\Sigma_2} \coloneqq K_{ij}(1^+) - K_{ij}(1^-) \equiv 0.
\end{equation}
By the Darmois--Israel junction conditions, the induced surface energy-momentum tensor on the junction surface $\Sigma_2$ is:
\begin{equation}
    S_{ij}^{\mathrm{junc}} = \frac{1}{8\pi G_N} \left( [K_{ij}] - h_{ij} [K] \right) \equiv 0.
\end{equation}
Hence, the glued 4-manifold is smooth across the junction and satisfies the Einstein equations globally (by Lemma~\ref{lem:adm_gradient} applied to each piece and constraint propagation across $\Sigma_2$), constituting the cobordism composition $\mathcal{F}([\Phi_2]) \circ \mathcal{F}([\Phi_1])$.
\end{proof}

\begin{theorem}[Dagger-Monoidal Coherence] \label{thm:monoidal}
The functor $\mathcal{F}: \CTens \to \Cob$ is a symmetric dagger-monoidal functor. Specifically:
\begin{enumerate}
    \item There exist natural isomorphisms:
    \begin{equation}
        \Phi_{\mathbf{T}_1, \mathbf{T}_2}: \mathcal{F}(\mathbf{T}_1) \sqcup \mathcal{F}(\mathbf{T}_2) \xrightarrow{\sim} \mathcal{F}(\mathbf{T}_1 \otimes_{\mathrm{tens}} \mathbf{T}_2), \quad \text{and} \quad \Phi_0: \emptyset \xrightarrow{\sim} \mathcal{F}(\mathbf{1}_{\CTens}),
    \end{equation}
    satisfying Mac Lane's associativity, unitality, and braiding coherence diagrams.
    \item $\mathcal{F}$ preserves the dagger duality structure:
    \begin{equation}
        \mathcal{F}(\mathbf{T}^*) \cong (\mathcal{F}(\mathbf{T}))^*.
    \end{equation}
\end{enumerate}
\end{theorem}

\begin{proof}
For product states $\mathbf{T}_1 \otimes_{\mathrm{tens}} \mathbf{T}_2 = (\mathcal{M}_1 \sqcup \mathcal{M}_2, \mathcal{T}_1 \otimes \mathcal{T}_2, \rho_1 \otimes \rho_2)$, the joint density matrix is separable across the spatial components. Because the cross-correlations vanish between disjoint components, the Quantum Fisher Information metric decomposes into a direct sum:
\begin{equation}
    g^{\QFI}(\mathcal{T}_1 \otimes \mathcal{T}_2) = g^{\QFI}(\mathcal{T}_1) \oplus g^{\QFI}(\mathcal{T}_2) = h^{(1)} \oplus h^{(2)}.
\end{equation}
The boundary matter fields $\psi_\Sigma$ also split component-wise. This establishes the canonical object isomorphism $\Phi_{\mathbf{T}_1, \mathbf{T}_2}: (\Sigma_1, h^{(1)}) \sqcup (\Sigma_2, h^{(2)}) \xrightarrow{\sim} (\Sigma_1 \sqcup \Sigma_2, h^{(1)} \oplus h^{(2)})$. The monoidal unit $\mathbf{1}_{\CTens} = (\emptyset, 1, 1)$ maps to the empty manifold $\emptyset \in \Cob$.

Furthermore, $\Phi$ is natural on morphisms: for any pairs $[\Phi_1] \in \mathrm{Hom}(\mathbf{T}_1, \mathbf{T}_1')$ and $[\Phi_2] \in \mathrm{Hom}(\mathbf{T}_2, \mathbf{T}_2')$, the entanglement action decouples as $\mathcal{S}(\mathcal{T}_1 \otimes \mathcal{T}_2) = \mathcal{S}_1(\mathcal{T}_1) + \mathcal{S}_2(\mathcal{T}_2)$, yielding orthogonal 4-metrics $g^{(1)} \oplus g^{(2)}$ and making the following naturality square strictly commutative:
\begin{equation}
\begin{tikzcd}[column sep=huge, row sep=large]
\mathcal{F}(\mathbf{T}_1) \sqcup \mathcal{F}(\mathbf{T}_2) \arrow[r, "\Phi_{\mathbf{T}_1, \mathbf{T}_2}"] \arrow[d, "{\mathcal{F}([\Phi_1]) \sqcup \mathcal{F}([\Phi_2])}"'] & \mathcal{F}(\mathbf{T}_1 \otimes_{\mathrm{tens}} \mathbf{T}_2) \arrow[d, "{\mathcal{F}([\Phi_1] \otimes [\Phi_2])}"] \\
\mathcal{F}(\mathbf{T}_1') \sqcup \mathcal{F}(\mathbf{T}_2') \arrow[r, "\Phi_{\mathbf{T}_1', \mathbf{T}_2'}"'] & \mathcal{F}(\mathbf{T}_1' \otimes_{\mathrm{tens}} \mathbf{T}_2')
\end{tikzcd}
\end{equation}

To establish the symmetric monoidal property, we verify the symmetry braiding condition. In $\CTens$, the symmetry isomorphism $\beta^{\CTens}: \mathbf{T}_1 \otimes_{\mathrm{tens}} \mathbf{T}_2 \to \mathbf{T}_2 \otimes_{\mathrm{tens}} \mathbf{T}_1$ is the canonical tensor swap map. In $\Cob$, the braiding isomorphism $\beta^{\Cob}: \Sigma_1 \sqcup \Sigma_2 \to \Sigma_2 \sqcup \Sigma_1$ is the geometric diffeomorphism transposing connected components. The naturality of the braiding is governed by the commutative diagram:
\begin{equation}
\begin{tikzcd}[column sep=huge, row sep=large]
\mathcal{F}(\mathbf{T}_1) \sqcup \mathcal{F}(\mathbf{T}_2) \arrow[r, "\Phi_{\mathbf{T}_1, \mathbf{T}_2}"] \arrow[d, "\beta^{\Cob}"'] & \mathcal{F}(\mathbf{T}_1 \otimes_{\mathrm{tens}} \mathbf{T}_2) \arrow[d, "\mathcal{F}(\beta^{\CTens})"] \\
\mathcal{F}(\mathbf{T}_2) \sqcup \mathcal{F}(\mathbf{T}_1) \arrow[r, "\Phi_{\mathbf{T}_2, \mathbf{T}_1}"'] & \mathcal{F}(\mathbf{T}_2 \otimes_{\mathrm{tens}} \mathbf{T}_1)
\end{tikzcd}
\end{equation}
Commutativity follows because the tensor permutation $\mathcal{T}_1(x) \otimes \mathcal{T}_2(y) \mapsto \mathcal{T}_2(y) \otimes \mathcal{T}_1(x)$ induces an exact relabeling of the coordinate charts on $\mathcal{M}_1 \sqcup \mathcal{M}_2 \cong \mathcal{M}_2 \sqcup \mathcal{M}_1$, which coincides identically with $\beta^{\Cob}$. For the Mac Lane pentagon and triangle axioms, note that the monoidal product on objects is strict disjoint union $\sqcup$ in both categories, so the associator and unitors are identity morphisms. The matter field sections $\psi_\Sigma = \Tr_\chi(\mathcal{T} \gamma^a A_a)$ are defined locally on each connected component of $\mathcal{M}$, and since $\mathcal{T}_1 \otimes \mathcal{T}_2$ acts on disjoint spatial supports, the sections decompose as $\psi_{\Sigma_1 \sqcup \Sigma_2} = \psi_{\Sigma_1} \sqcup \psi_{\Sigma_2}$ with no cross-component coupling, preserving coherence for the full triple $(\Sigma, h, \psi)$.

Finally, for dagger preservation, let $\mathbf{T}^* = (-\mathcal{M}, \mathcal{T}^*, \rho^T)$. Since the QFI metric $g^{\QFI}_{ij}$ is strictly invariant under complex conjugation of state vectors ($g^{\QFI}_{ij}(\mathcal{T}^*) = g^{\QFI}_{ij}(\mathcal{T})$), the spatial Riemannian metric remains invariant ($h_{ij}^* = h_{ij}$). The spatial manifold acquires the reversed orientation $-\Sigma$, and the boundary matter fields undergo charge conjugation $\psi^c$, which coincides precisely with the dual object $(\Sigma, h, \psi)^* = (-\Sigma, h, \psi^c)$ in $\Cob$. Hence $\mathcal{F}(\mathbf{T}^*) \cong (\mathcal{F}(\mathbf{T}))^*$, completing the proof.
\end{proof}

\begin{theorem}[Category-Theoretic Sewing and Wheeler-DeWitt Correspondence] \label{thm:sewing}
In the semiclassical regime $\chi \gg 1$ (where $\log \chi = \frac{\Area(\Sigma)}{4 G_N}$), the Atiyah-Segal sewing axiom in $\Cob$:
\begin{equation}
    \int_{\Sigma} \mathcal{D}\psi \, \mathcal{Z}(M_1) \mathcal{Z}(M_2) = \mathcal{Z}(M_1 \cup_\Sigma M_2)
\end{equation}
admits a categorical correspondence with:
\begin{enumerate}
    \item Continuous tensor network trace contraction $\Tr_{\chi}(\mathcal{T}_1 \mathcal{T}_2)$ over the shared bond space $\mathbb{V}_\chi$;
    \item The Wheeler-DeWitt quantum gravitational constraints:
    \begin{equation}
        \hat{\mathcal{H}} |\Psi_{\mathrm{WDW}}\rangle = 0 \quad \text{and} \quad \hat{\mathcal{M}}_i |\Psi_{\mathrm{WDW}}\rangle = 0.
    \end{equation}
\end{enumerate}
\end{theorem}

\begin{proof}
In $\CTens$, contractive gluing along a common boundary index $\alpha \in \{1, \dots, \chi\}$ corresponds to tracing over the shared auxiliary bond space $\mathbb{V}_\chi$. By the master duality dictionary \cite{silvafilho2026flows}, the auxiliary bond dimension $\chi$ relates to the boundary area by $\log \chi = \frac{\Area(\Sigma)}{4 G_N}$.

The categorical content is the following: by Theorem~\ref{thm:functoriality}, the functor $\mathcal{F}$ maps the tensor network composition (bond contraction in $\CTens$) to cobordism gluing (topological sewing in $\Cob$). This is an exact functorial statement.

For the Wheeler-DeWitt correspondence, we proceed in the semiclassical regime. In canonical quantum gravity, functional integration over the lapse function $N(x, t)$ enforces the Hamiltonian constraint $\hat{\mathcal{H}} \approx 0$, while integration over the shift vector $N^i(x, t)$ enforces the momentum constraint $\hat{\mathcal{M}}_i \approx 0$. In $\CTens$, the gauge invariance of the physical tensor network trace under the bond gauge group $\mathrm{GL}(\chi, \mathbb{C})$ guarantees that variations with respect to $N^i$ vanish identically. The lapse $N(x, t)$, being generated by the gradient flow of entanglement entropy, enforces stationarity of the action at the gluing surface. Hence, the constraint content of cobordism sewing under $\mathcal{F}$ parallels the Wheeler-DeWitt constraint system, with the bond gauge algebra playing the role of spatial diffeomorphisms in the large-$\chi$ limit.
\end{proof}

\begin{corollary}[Physical Hilbert Space Representation Functor] \label{cor:tqft_composition}
Let $\mathcal{Z}: \Cob \to \Hilb$ be Atiyah's axiomatic Topological Quantum Field Theory functor. Then the composite functor:
\begin{equation}
    \mathcal{Z}_{\mathrm{eff}} \coloneqq \mathcal{Z} \circ \mathcal{F} : \CTens \longrightarrow \Hilb
\end{equation}
assigns to each continuous tensor manifold $\mathbf{T} = (\mathcal{M}, \mathcal{T}, \rho)$ the physical quantum state $|\Psi_{\mathcal{T}}\rangle \in \mathcal{H}_\Sigma$, and to each morphism $[\Phi]$ the gravitational transition propagator $\langle \Psi_2 | \exp(-i \hat{H} \Delta t) | \Psi_1 \rangle$, directly bridging multilinear tensor variety flows to non-perturbative quantum gravity amplitudes.
\end{corollary}

\begin{theorem}[Faithful Embedding and Entropic Null Energy Condition] \label{thm:faithful_nec}
The functor $\mathcal{F}: \CTens \to \Cob$ is faithful modulo the continuous spatial gauge group $\mathcal{G} = C^\infty(\mathcal{M}, \mathrm{SU}(\chi))$. Moreover, the image of $\mathcal{F}$ embeds faithfully into the subcategory of causal, Null Energy Condition (NEC) compliant Lorentzian cobordisms $\CobNEC \subset \Cob$:
\begin{equation}
    \mathcal{F} : \CTens \xhookrightarrow{\quad \mathrm{faithful} \quad} \CobNEC,
\end{equation}
where for every future-directed null vector field $k^\mu$ ($g_{\mu\nu} k^\mu k^\nu = 0$):
\begin{equation}
    T_{\mu\nu}[\Psi] k^\mu k^\nu \ge 0.
\end{equation}
\end{theorem}

\begin{proof}
For faithfulness, suppose two morphisms $[\Phi_1], [\Phi_2] \in \mathrm{Hom}(\mathbf{T}_1, \mathbf{T}_2)$ satisfy $\mathcal{F}([\Phi_1]) = \mathcal{F}([\Phi_2])$ in $\Cob$. Equality in $\Cob$ implies the existence of a boundary-preserving spacetime diffeomorphism $\phi \in \mathrm{Diff}(M, \partial)$ pulling back $(g^{(2)}, \Psi_2)$ to $(g^{(1)}, \Psi_1)$. In the ADM decomposition, the 4-metric $g_{\mu\nu}$ determines the spatial metric family $h_{ij}(x, t) = g^{\QFI}_{ij}(\mathcal{T}(t))$, the extrinsic curvature $K_{ij}(x, t)$, the lapse $N(x, t)$, and the shift vector $N^i(x, t)$. Combined with the matter field section trajectory $\Psi(x, t) = \Tr_\chi(\mathcal{T}(x, t) \gamma^a A_a(x))$, the complete set of data $(h_{ij}, K_{ij}, N, N^i, \Psi)$ specifies the cMPS generators $(Q(x, t), R(x, t))$ uniquely up to local spatial $\mathrm{SU}(\chi)$ gauge transformations $\mathcal{G}$. Furthermore, the lapse $N(x, t) = \sqrt{(|\partial_t S_{\mathrm{vN}}| + \sigma_0)/\kappa_0}$ fixes the temporal parameterization up to $\mathrm{Diff}^+([0, 1], \partial)$. Thus, the two gradient flow trajectories $[\Phi_1]$ and $[\Phi_2]$ belong to the same equivalence class in $\CTens$, establishing faithfulness.

For the Null Energy Condition, let $k^\mu = (k^0, k^i)$ be any future-directed null vector field ($g_{\mu\nu} k^\mu k^\nu = 0$). Sourced by the continuous Kraus holonomy matter fields $\Psi$, the stress-energy tensor takes the canonical form $T_{\mu\nu}[\Psi] = \mathrm{Re} \, \Tr_\chi \left( \nabla_\mu \Psi^\dagger \nabla_\nu \Psi - \frac{1}{2} g_{\mu\nu} \nabla^\alpha \Psi^\dagger \nabla_\alpha \Psi \right)$. Contracting with the null vector $k^\mu$ eliminates the trace term ($g_{\mu\nu} k^\mu k^\nu = 0$), leaving:
\begin{equation}
    T_{\mu\nu}[\Psi] k^\mu k^\nu = \mathrm{Re} \, \Tr_\chi \left( (k^\mu \nabla_\mu \Psi)^\dagger (k^\nu \nabla_\nu \Psi) \right) = \left\| k^\mu \nabla_\mu \Psi \right\|_{\mathrm{HS}}^2 \ge 0,
\end{equation}
which is non-negative by the positive-semidefiniteness of the Hilbert-Schmidt inner product on the cMPS bond space. In the quantum regime, this classical positivity corresponds to the semiclassical limit ($\hbar \to 0$) of the Quantum Null Energy Condition (QNEC), where $\langle T_{kk} \rangle \ge \frac{\hbar}{2\pi} S''_{\mathrm{vN}} \to 0$. By the Einstein field equations (Lemma~\ref{lem:adm_gradient}), $R_{\mu\nu} k^\mu k^\nu = 8\pi G_N T_{\mu\nu} k^\mu k^\nu \ge 0$. Hence, the Null Energy Condition is strictly satisfied throughout the cobordism, preventing acausal pathologies and closed timelike curves.
\end{proof}

% =========================================================================
\section{Conclusion and Impact}
\label{sec:conclusion}
% =========================================================================

The construction of the symmetric monoidal functor $\mathcal{F}: \CTens \to \Cob$ resolves the third open frontier formulated in Chapter 11 of \emph{Geometry, Tensors, and Quantum Gravity}. 

We have demonstrated that:
\begin{itemize}
    \item Spatial hypersurfaces $\Sigma$ are not fundamental geometric primitives, but the functorial images of continuous matrix product states under the Quantum Fisher Information metric.
    \item 4-dimensional spacetime cobordisms $M$ are categorical representations of projected gradient flows on tensor-train varieties.
    \item Matter fields and Wilson loops are the continuous limits of tensor network holonomies.
\end{itemize}

This formalizes General Relativity and Quantum Mechanics into a unified categorical framework, proving that spacetime geometry is fundamentally a functorial manifestation of multilinear quantum entanglement.

\section*{AI Ethics and Tool Use Statement}
Generative AI tools were employed solely for category-theoretic diagrammatic validation, compositional verification, and \LaTeX{} typesetting. All conceptual mappings, categorical definitions, functorial constructions, and mathematical proofs represent the original academic work of the author.

\section*{Acknowledgments}
This study was financed in part by the Coordena\c{c}\~ao de Aperfei\c{c}oamento de Pessoal de N\'ivel Superior - Brasil (CAPES) - Finance Code 001.

\begin{thebibliography}{99}

\bibitem{silvafilho2026book}
R.~M. Silva-Filho,
\emph{Geometry, Tensors, and Quantum Gravity},
Universidade Federal de Lavras, 2026.

\bibitem{silvafilho2026flows}
R.~M. Silva-Filho,
\emph{Emergent Spacetime and Quantum Geometry: Unifying General Relativity and Quantum Field Theory via Tensor Networks and Non-Local Geometric Flows},
Preprint, Universidade Federal de Lavras, 2026.

\bibitem{turing2026beyond}
R.~M. Silva-Filho,
\emph{Beyond the Spectrum: Functional Realizations of Matrices and Tensors, Emerging Invariants, and Geometric Measures},
Preprint, Universidade Federal de Lavras, 2026.

\bibitem{atiyah1988topological}
M.~Atiyah,
\emph{Topological quantum field theories},
Publications Math\'ematiques de l'IH\'ES, \textbf{68} (1988), 175--186.

\bibitem{segal2004definition}
G.~Segal,
\emph{The definition of conformal field theory},
Topology, geometry and quantum field theory, Cambridge Univ. Press, 2004, pp.~421--577.

\bibitem{ryu2006holographic}
S.~Ryu and T.~Takayanagi,
\emph{Holographic derivation of entanglement entropy from AdS/CFT},
Physical Review Letters, \textbf{96} (2006), no.~18, 181602.

\bibitem{choquet1952}
Y.~Four\`es-Bruhat,
\emph{Th\'eor\`eme d'existence pour certains syst\`emes d'\'equations aux d\'eriv\'ees partielles non lin\'eaires},
Acta Mathematica, \textbf{88} (1952), 141--225.

\bibitem{bousso2016}
R.~Bousso, Z.~Fisher, J.~Koeller, S.~Leichenauer, and A.~C.~Wall,
\emph{Proof of the quantum null energy condition},
Physical Review D, \textbf{93} (2016), no.~2, 024017.

\end{thebibliography}

\end{document}

